
\section{Annotation Pipeline Details}
\label{app:annotation}

\begin{figure}[h!]
  \centering
  \includegraphics[width=\linewidth]{figure2.001.jpeg}
  \caption{
    \textbf{Three-phase annotation pipeline.}
    Phase 1 (SAM-assisted manual labeling) trains an initial model; Phase 2
    (model-in-the-loop) accelerates throughput $\sim$4$\times$; Phase 3
    (retroactive quality control) reconciles labels against final model predictions.
    Each phase's output seeds the next, jointly solving scale and label quality.
    Inter-annotator agreement statistics and adjudication rules are in
    Appendix~\ref{app:annotation}.
  }
  \label{fig:annotation}
\end{figure}

This appendix documents the three-phase pipeline summarized in
Section~\ref{sec:annotation}: the per-phase tooling
(Figure~\ref{fig:annotation}), the Phase~3 adjudication rules that
reconcile human labels with model predictions
(Table~\ref{tab:adjudication}), and the gold-image protocol and
Cohen's $\kappa$ trajectory that monitor annotator drift across batches
(Table~\ref{tab:kappa}). Together these mechanisms support the
human-verified label claim made against S-EO~\cite{seo} in
Table~\ref{tab:dataset_comparison}.

\subsection{Three-Phase Pipeline}
\label{app:phases}

Seven trained annotators produced shadow masks in CVAT across
$45$ sessions of $105$ images each ($100$ working images $+$ $5$
gold images with frozen reference masks; \S\ref{app:gold}). Sessions
proceed in three sequential phases (Figure~\ref{fig:annotation}); each
phase's output seeds the next.

\textbf{Phase 1: SAM-assisted manual labeling.} Annotators draw shadow
masks on raw $384{\times}384$ patches. The CVAT AI Interactor \cite{cvat}, powered
by the Segment Anything Model (SAM)~\cite{sam}, proposes a
candidate boundary from a single annotator click; the annotator
accepts, refines, or rejects it. Manual polygon drawing is available
for thin-linear and irregular shadows where SAM under-segments. Each
segment receives a shadow-type label from the six-category taxonomy
(\S\ref{app:taxonomy}). Phase~1 spans the first $10$ sessions and
is restricted to Chicago and Miami imagery balanced across the two
native resolutions; Phoenix is admitted only after the label
distribution stabilizes (gold-image $\kappa \geq 0.7$ for a
majority of annotators). Mean throughput is 4.8 minutes per
image after a one-session per-annotator ramp-up.

\textbf{Phase 2: Model-in-the-loop acceleration.} After Phase~1, all
labels train an interim shadow detector. The detector is one of
\{MAMNet, OGLANet, DINOv3\}, drawn at random for each new $100$-image
batch; rotating architectures limits annotator entrainment to any
single model's failure modes. Predictions from the interim model are
loaded into CVAT as the starting layer; annotators delete false
positives, refine boundaries, and add missed segments rather than
redraw from scratch. The interim model is retrained at session
boundaries on the cumulative labeled corpus. Throughput rises to
$75$ seconds per image, a $\sim\!4\times$ acceleration over
Phase~1.

\textbf{Phase 3: Retroactive quality control.} After all working
sessions complete, a final shadow detector is trained on the full
Phase~1$+$2 corpus and used to revisit every annotated segment. For
each image we compute the per-segment Intersection-over-Union (IoU)
between each human annotation $j$ and each model prediction $i$. Each
(annotation, prediction) pair receives one of four statuses
(Table~\ref{tab:adjudication}). Statuses are color-coded in CVAT so
annotators triage at a glance: green segments are skipped; yellow,
blue, and red segments are routed to a second annotator pass for
adjudication. The status taxonomy is released with the dataset for
downstream provenance analyses; for the LOCO benchmark
(\S\ref{sec:loco}), all final segments are merged into the binary
foreground class regardless of their adjudication route.

\begin{table}[h]
  \centering
  \small
  \setlength{\tabcolsep}{4pt}
  \caption{
    \textbf{Phase 3 adjudication rules.}
    Each (annotation, prediction) pair on an image is assigned one of
    four statuses based on their IoU and overlap, and routed to the
    corresponding annotator action.
    $j$: human annotation; $i$: model prediction.
    Distribution column reports the share of all Phase 3 segments
    falling into each status.
  }
  \label{tab:adjudication}
  \begin{tabular}{@{}lp{3.1cm}p{2cm}p{3.4cm}c@{}}
    \toprule
    \textbf{Status (color)} & \textbf{Trigger} & \textbf{Kept segment} & \textbf{Annotator action} & \textbf{Share} \\
    \midrule
    Correct (green) &
      $j \cap i \neq \emptyset$, IoU\,$>0.5$ &
      $j$ &
      Verify, usually skip &
      $65\%$ \\
    Potential\_Annotated (yellow) &
      $j \cap i \neq \emptyset$, IoU\,$\leq 0.5$ &
      $j$ (edit shape) &
      Adjust boundary &
      $21\%$ \\
    Potential\_Model (blue) &
      $i$ has no overlapping $j$ &
      $i$ if validated &
      Confirm or reject; assign type &
      $9\%$ \\
    Incorrect (red) &
      $j$ has no overlapping $i$ &
      $j$ if validated &
      Verify segment is shadow &
      $5\%$ \\
    \bottomrule
  \end{tabular}
\end{table}

\subsection{Gold-Set Protocol and Inter-Annotator Agreement}
\label{app:gold}

\textbf{Gold images.} Every $100$-image working batch is paired with
$5$ gold images that carry frozen reference masks. Gold images are
drawn from previously closed batches that have already cleared
Phase~3, and their reference masks are independently re-adjudicated
by two senior annotators before the gold image enters circulation.
Once frozen, a gold mask is not modified. Gold images are interleaved
with working images and are not visually distinguishable from them in
CVAT, so annotator behavior on gold matches their behavior on the rest
of the batch.

\textbf{Cohen's $\kappa$ trajectory.} For each annotator and each
batch, we compute Cohen's $\kappa$ on binary shadow/background
agreement between the annotator's mask and the frozen gold reference,
pooled over the $5$ gold images. Mean $\kappa$ across the seven
annotators rises from Phase~1 to Phase~2 as model seeding stabilizes
the label distribution, and rises further after Phase~3 adjudication
(Table~\ref{tab:kappa}). All Phase~2 and Phase~3 values fall within the
substantial-to-almost-perfect range
($\kappa\!\geq\!0.80$)~\cite{landis1977measurement}.

\begin{table}[h]
  \centering
  \small
  \setlength{\tabcolsep}{6pt}
  \caption{
    \textbf{Inter-annotator agreement across phases.}
    Mean Cohen's $\kappa$ on gold images across the seven annotators,
    with the across-annotator range and the per-batch first-submission
    acceptance rate (\S\ref{app:gates}). Phase~3 reports
    post-adjudication agreement against the final reconciled corpus.
  }
  \label{tab:kappa}
  \begin{tabular}{@{}lcccc@{}}
    \toprule
    \textbf{Phase} & \textbf{Sessions} & \textbf{Mean $\kappa$} & \textbf{Range across annotators} & \textbf{First-submission accept} \\
    \midrule
    Phase 1 (manual)            & $1\text{--}10$  & $0.71$ & $[0.62,\,0.78]$ & $84\%$ \\
    Phase 2 (model-in-the-loop) & $11\text{--}45$ & $0.84$ & $[0.79,\,0.89]$ & $96\%$ \\
    Phase 3 (post-adjudication) & ---                   & $0.89$ & $[0.85,\,0.92]$ & --- \\
    \bottomrule
  \end{tabular}
\end{table}

\textbf{Acceptance gates.}
\label{app:gates}
A batch is accepted when its mean gold mask-IoU against the frozen
reference is at least $0.90$; spot-checked on $5$--$10$ randomly
selected working images by a senior annotator before entering the
corpus. A batch with mean gold IoU between $0.80$ and $0.90$ is
conditionally accepted: the annotator receives written feedback on the
failed gold images, and the next batch dispatched to that annotator
includes $5$ additional gold images. A batch with mean gold IoU below
$0.80$ is rejected and the failed images returned for redo; a second
failed submission triggers a paired-review session with a senior
annotator before the next solo batch is dispatched. Across all
sessions, $96\%$ of batches were accepted on first submission, and
$4\%$ were resubmitted.

\subsection{Annotator Recruitment and Instructions}
\label{app:annotator-instructions}

Annotation was paid, on-demand work carried out in India. The full text of the recruitment and task instructions given to annotators follows; amounts are in Indian rupees (INR). In practice every accepted image was paid INR~20, and each review by a second annotator was paid a further INR~20.

\begin{quote}\small
\textbf{Mark Shadow Areas on Images (Web Tool)}

\textit{Why this work matters.} We're building a high-quality research dataset that helps improve mapping and city planning. Your careful work directly contributes to real-world research used by universities and labs.

\textit{What you'll do.} Open a small set of images (satellite view) in a web tool. Mark areas that are in shadow by drawing/adjusting shapes. Save and move to the next image. That's it. No coding. No prior mapping experience required.

\textit{Training \& support.} Short written guide with pictures. Quick sample demo. We provide fast feedback so you can improve and earn the quality bonus.

\textit{Payment (per accepted image).} INR 20 per image. Each session is ${\sim}100$ images, so the typical total is INR 1,500--2,000. Payments are made after review of each session.

\textit{How we review quality.} Each batch passed through a validation AI model. We do a short spot-check on a few others. If any images need fixes, we'll send them back with quick notes; you'll be paid for the final images once approved.

\textit{Selection process (short test).} You'll get a mini test (5 images) in the same web tool. If the test looks good, we'll invite you for paid work immediately. Top performers get more sessions each week.

\end{quote}